% Generisano: uređuj beleske/sNNN.tex i naslove u manifestu.
\documentclass[11pt,a4paper]{article}
% Zajednički preambulum za dokument za čitanje; nije Beamer tema.
\usepackage[a4paper,margin=20mm,headheight=15pt]{geometry}
\usepackage{lmodern,fontspec}
\usepackage[serbian]{babel}
\setmainfont{Latin Modern Sans}
\setsansfont{Latin Modern Sans}
\setmonofont{Latin Modern Mono}
\usepackage{amsmath,amssymb,mathtools,graphicx,xcolor}
\usepackage{array,booktabs,tabularx,multirow,makecell,listings,fancyhdr}
\pagestyle{fancy}\fancyhf{}
\newcommand{\BeleskeTekuciID}{NEOZNACEN}
\fancyhead[L]{\small Beleške uz slajd \texttt{\expandafter\detokenize\expandafter{\BeleskeTekuciID}}}
\fancyfoot[R]{\small\thepage}
\setlength{\parindent}{0pt}\setlength{\parskip}{6pt}
\newwrite\BeleskeMapa
\immediate\openout\BeleskeMapa=\jobname.notes-map.tsv
\AddToHook{shipout/before}{%
  \immediate\write\BeleskeMapa{\BeleskeTekuciID\space\thepage}}
\AddToHook{enddocument/afterlastpage}{\immediate\closeout\BeleskeMapa}
\newcommand{\BeleskeID}[1]{%
  \def\BeleskeZadatiID{#1}%
  \ifx\BeleskeTekuciID\BeleskeZadatiID\else
    \PackageError{beleske}{Pogresan ID beleske: #1}{Proveri mapiranje slajda.}%
  \fi}
\newcommand{\BeleskaSlajda}[4]{%
  \clearpage\gdef\BeleskeTekuciID{#1}%
  {\Large\bfseries Slajd \textnormal{\texttt{\detokenize{#1}}}: #3\par}\medskip
  \includegraphics[page=#2,width=\linewidth]{\BeleskePrezentacija}\par\medskip
  \input{#4}}

\newcommand{\BeleskePrezentacija}{build/04_brojni_sistemi.pdf}
\begin{document}
\BeleskaSlajda{04-s001}{1}{Brojni sistemi}{beleske/s001.tex}
\BeleskaSlajda{04-s002}{2}{Pozicioni brojni sistemi}{beleske/s002.tex}
\BeleskaSlajda{04-s003}{3}{Celobrojna i razlomljena vrednost}{beleske/s003.tex}
\BeleskaSlajda{04-s004}{4}{Granica razlomljene vrednosti}{beleske/s004.tex}
\BeleskaSlajda{04-s005}{5}{Prelazak u decimalni sistem}{beleske/s005.tex}
\BeleskaSlajda{04-s006}{6}{Primer: iz osnove 7 u osnovu 10}{beleske/s006.tex}
\BeleskaSlajda{04-s007}{7}{Prelazak iz decimalnog sistema}{beleske/s007.tex}
\BeleskaSlajda{04-s008}{8}{Celobrojni deo: uzastopno deljenje}{beleske/s008.tex}
\BeleskaSlajda{04-s009}{9}{Da li je celobrojni zapis uvek konačan?}{beleske/s009.tex}
\BeleskaSlajda{04-s010}{10}{Razlomljeni deo: uzastopno množenje}{beleske/s010.tex}
\BeleskaSlajda{04-s011}{11}{Da li je razlomljeni zapis uvek konačan?}{beleske/s011.tex}
\BeleskaSlajda{04-s012}{12}{Uslov konačnog razlomljenog zapisa}{beleske/s012.tex}
\BeleskaSlajda{04-s013}{13}{Primer: decimalni broj u binarnom sistemu}{beleske/s013.tex}
\BeleskaSlajda{04-s014}{14}{Deljenje u osnovi 7}{beleske/s014.tex}
\BeleskaSlajda{04-s015}{15}{Primer: uzastopno deljenje u osnovi 7}{beleske/s015.tex}
\BeleskaSlajda{04-s016}{16}{Kompatibilni brojni sistemi}{beleske/s016.tex}
\BeleskaSlajda{04-s017}{17}{Binarni, oktalni i heksadecimalni sistem}{beleske/s017.tex}
\BeleskaSlajda{04-s018}{18}{Cifre kompatibilnih sistema}{beleske/s018.tex}
\BeleskaSlajda{04-s019}{19}{Grupisanje celobrojnih cifara}{beleske/s019.tex}
\BeleskaSlajda{04-s020}{20}{Grupisanje razlomljenih cifara}{beleske/s020.tex}
\BeleskaSlajda{04-s021}{21}{Prikaz brojeva u digitalnom sistemu}{beleske/s021.tex}
\BeleskaSlajda{04-s022}{22}{Konačan broj informacija}{beleske/s022.tex}
\BeleskaSlajda{04-s023}{23}{Bit, bajt i reč}{beleske/s023.tex}
\BeleskaSlajda{04-s024}{24}{Koje informacije se nalaze u registru?}{beleske/s024.tex}
\BeleskaSlajda{04-s025}{25}{Registar kao binarni broj}{beleske/s025.tex}
\BeleskaSlajda{04-s026}{26}{Gde je binarna tačka?}{beleske/s026.tex}
\BeleskaSlajda{04-s027}{27}{Digresija: skaliranje vrednosti registra}{beleske/s027.tex}
\BeleskaSlajda{04-s028}{28}{Negativne vrednosti brojeva}{beleske/s028.tex}
\BeleskaSlajda{04-s029}{29}{Znak i apsolutna vrednost}{beleske/s029.tex}
\BeleskaSlajda{04-s030}{30}{Prikaz sa ofsetom}{beleske/s030.tex}
\BeleskaSlajda{04-s031}{31}{Prvi komplement}{beleske/s031.tex}
\BeleskaSlajda{04-s032}{32}{Drugi komplement}{beleske/s032.tex}
\BeleskaSlajda{04-s033}{33}{Komplement maksimalne vrednosti}{beleske/s033.tex}
\BeleskaSlajda{04-s034}{34}{Komplement osnove}{beleske/s034.tex}
\BeleskaSlajda{04-s035}{35}{Ponovljeno komplementiranje}{beleske/s035.tex}
\BeleskaSlajda{04-s036}{36}{Težina bita znaka u drugom komplementu}{beleske/s036.tex}
\BeleskaSlajda{04-s037}{37}{Drugi komplement i prikaz sa ofsetom}{beleske/s037.tex}
\BeleskaSlajda{04-s038}{38}{Negativna vrednost sa fiksnom tačkom}{beleske/s038.tex}
\BeleskaSlajda{04-s039}{39}{Prikaz sa pokretnom tačkom}{beleske/s039.tex}
\end{document}
